% LaTex source file for LAA-D-14-00116 
%
\documentclass [review] {elsarticle}
\usepackage{latexsym,amsmath,amssymb,graphicx,slashbox}
\usepackage{color,url}
\def\bea#1{\begin{array}{#1}}
\def\ea{\end{array}}
\def\reff#1{(\ref{#1})}
\def\T{^\top}
\def\eps{\varepsilon}
\def\R{{\mathbb R}}
\def\ve#1{{\rm vec }(#1)}
\newtheorem{defn}{Definition}[section]
\newtheorem{thm}{Theorem}[section]
\newtheorem{lem}{Lemma}[section]
\newtheorem{cly}{Corollary}[section]
\newtheorem{rem}{Remark}[section]
\newtheorem{prop}{Proposition}[section]
\newtheorem{obs}{Observation}[section]
\newtheorem{alg}{Algorithm}[section]
\newproof{proof}{Proof}
\def\qf#1#2{#1\T  #2 #1}
\def\ker{\mbox{ker }}
\def\sgn{\mbox{sgn\,}}
\def\argmin{\mbox{argmin }}
\def\relint{\mbox{relint }}
\def\Diag{\mbox{Diag }}
\def\diag{\mbox{diag }}
\def\det{\mbox{det }}
\def\conv{\mbox{conv\,}}
\def\tr{\mbox{trace }}
\def\rank{\mbox{rank }}
\def\ext{\mbox{ext }}
\def\beq#1{\begin{equation}\label{#1}}
\def\eeq{\end{equation}}
\def\bep{\noindent {\sc Proof. }}
\def\ep{\hfill$\Box$}
\def\che{\textstyle{\frac 12}}
\def\lk{\left\{}
\def\rk{\right \} }
\def\du#1#2{\langle {#1} , {#2} \rangle}
\def\norm#1{ \| #1 \|}
\def\mr#1{(\ref{#1})}
\def\xra#1{\xrightarrow{#1}}
\def\VV{{\mathcal V}}
\def\cs#1{{{\cal C}^{#1*}}}
\def\cc#1{{\cal C}^{#1}}
\def\ct{{\cal T}}
\def\ce{{\cal E}}
\def\ca{{\cal A}}
\def\csn{{{\cal S}^n}}
\def\csnp{{{\cal S}^n_+}}
\def\x{\mathbf x}
\def\y{\mathbf y}
\def\z{\mathbf z}
\def\u{\mathbf u}
\def\v{\mathbf v}
\def\d{\mathbf d}
\def\w{\mathbf w}
\def\pp{\mathbf p}
\def\f{\mathbf f}
\def\o{\mathbf o}
\def\e{\mathbf e}
\def\b{\mathbf b}
\def\c{\mathbf c}
\def\a{{\mathbf a}}
\def\g{{\mathbf g}}
\def\n{{\mathbf n}}
\def\m{{\mathbf m}}
\def\q{\mathbf q}
\def\s{{\mathbf s}}
\def\t{{\mathbf t}}
\def\ot{\otimes}
\def\lm{\emptyset}
\def\Ab{{\mathsf A}}
\def\Bb{{\mathsf B}}
\def\Cb{{\mathsf C}}
\def\Db{{\mathsf D}}
\def\Eb{{\mathsf E}}
\def\Fb{{\mathsf F}}
\def\Gb{{\mathsf G}}
\def\Hb{{\mathsf H}}
\def\Ib{{\mathsf I}}
\def\Jb{{\mathsf J}}
\def\Kb{{\mathsf K}}
\def\Lb{{\mathsf L}}
\def\Mb{{\mathsf M}}
\def\Nb{{\mathsf N}}
\def\Ob{{\mathsf O}}
\def\Pb{{\mathsf P}}
\def\Qb{{\mathsf Q}}
\def\Rb{{\mathsf R}}
\def\Sb{{\mathsf S}}
\def\Tb{{\mathsf T}}
\def\Ub{{\mathsf U}}
\def\Vb{{\mathsf V}}
\def\Wb{{\mathsf W}}
\def\Xb{{\mathsf X}}
\def\Yb{{\mathsf Y}}
\def\Zb{{\mathsf Z}}
\def\inn#1{{#1}^\circ}
\def\cone{\mbox{cone }}
\def\supp#1{I(#1)}
\def\cpr{\mbox{cpr }}

\journal{Linear Algebra and its Applications}

\begin{document}

\begin{frontmatter}

\title{From seven to eleven:\\ Completely positive matrices with high cp-rank}

\author[univie]{Immanuel M. Bomze}


\author[univie]{Werner Schachinger}


\author[univie]{Reinhard Ullrich\corref{cor1}}

\ead{reinhard.ullrich@univie.ac.at}

\cortext[cor1]{Corresponding author}

\address[univie]{University of Vienna, Department of Statistics and Operations Research, 1090 Vienna, Austria}

\begin{abstract}
\normalsize
We study $n\times n$ completely positive matrices $\Mb$ on the boundary of the completely positive cone, namely those orthogonal to a copositive matrix $\Sb $ which generates a quadratic form with finitely many zeroes in the standard simplex. Constructing particular instances of $\Sb $, we are able to construct counterexamples to the famous Drew-Johnson-Loewy conjecture (1994) for matrices of order seven through eleven.
\end{abstract}


\begin{keyword}
copositive optimization \sep completely positive matrices \sep cp-rank \sep nonnegative factorization \sep circular symmetry

\MSC[2010] 15B48 \sep 90C25 \sep 15A23
\end{keyword}

\end{frontmatter}
 


\section{Introduction}

In this article we consider completely positive matrices $\Mb $ and their cp-rank.
An $n\times n$ matrix $\Mb $ is said to be {\em completely positive} if there exists a nonnegative  (not necessarily square) matrix $\Vb$ such that $\Mb =\Vb\Vb\T$.
Typically, a completely positive matrix $\Mb $ may have many such factorizations, and
the {\em cp-rank} of $\Mb $, $\cpr \Mb $, is the minimum number of columns in such a nonnegative factor $\Vb$ (for completeness, we define $\cpr \Mb  = 0$ if $\Mb $ is a square zero matrix and $\cpr \Mb  = \infty$ if $\Mb $ is not completely positive).
Completely positive matrices form a cone dual to the cone of
{\em copositive matrices}. An $n\times n$ matrix $\Sb$ is said to be copositive if $\qf{\x} \Sb  \ge 0$ for every nonnegative vector $\x\in \R^n_+$. Both cones are central in the rapidly evolving field of {\em copositive optimization} which links discrete and continuous optimization, and has
numerous real-world applications. For recent surveys and structured bibliographies, we refer to~\cite{Bomz11a,Bomz11b,Bure11,Duer10}, and for a fundamental text book to~\cite{Berm03}.

Determining the maximum possible cp-rank of $n\times n$ completely positive matrices,
$$p_n := \max\lk \cpr \Mb  : \Mb  \mbox{ is a completely positive }n\times n\mbox{ matrix}\rk ,$$
is still an open problem for general $n$. It is known \cite[Theorem~3.3]{Berm03} that
$p_n=n$ if $n\le 4$,  whereas this equality does no longer hold for $n\ge 5$.
Let $d_n := \left\lfloor \frac{n^2}4\right\rfloor$ and $s_n := {n+1\choose 2}-4$. For $n\ge 5$, it is known that~\cite{Shak13b}
\beq{pnbracket} d_n\le p_n \le s_n \, ,\eeq
and that $d_n=p_n$ in case $n=5$~\cite{Shak13a}.
It is still unknown whether $d_6=p_6$ although the bracket~\reff{pnbracket} was reduced in the recent paper~\cite{Shak13b} where also the upper bound $p_n\le s_n$ was established for the first time.


The famous Drew-Johnson-Loewy (DJL) conjecture~\cite{Drew94} is by now twenty years old. It states that $d_n=p_n$ is true for all $n\ge 5$, and some evidence in support of the DJL conjecture is found in~\cite{Berm98,Drew96a,Drew94,Loew03}, see also \cite[Section 3.3]{Berm03}. However, we will show in this paper that the DJL conjecture does not hold for $n\in \lk 7,8,9,10,11\rk$ by constructing examples which establish $p_n > d_n$.


The paper is organized as follows: In Section 2 we look at copositive matrices $\Sb $ which allow for finitely many (but many) zeroes $\q_i$ of the quadratic form $\qf\x \Sb $ over the standard simplex. Such matrices $\Sb $ lie on the boundary of the copositive cone, and elementary conic duality therefore tells us that there are nontrivial completely positive matrices $\Mb $ such that $\Mb \perp \Sb $ in the Frobenius inner product sense, and we will study the cp-rank of these $\Mb $. Section 3 deals with a particular construction of above mentioned copositive matrices $\Sb $ (they will be cyclically symmetric) in a way that many $\q_i$ can coexist, and in Section 4 we  {present} the second main result --  counterexamples to the DJL conjecture for  $7\le n\le 11$.
Let us mention here that such a counterexample for $n=7$ with cp-rank 14 was announced in 2002,  according to \cite[p.177]{Berm03}. The matrix there (which never got public) should have rank 5; by contrast, our matrix $\Mb $ in Example~1 will have full rank $7$, but also $\cpr \Mb  =14$ by mere coincidence.



Some notation and terminology:
we abbreviate $[r\!:\!s]=\lk r,r+1,\ldots ,s\rk$ for integers $r\le s$, and by $| S|$ the number of elements of a finite set $S$. For a function $f:T\to \R$ we abbreviate
$$\textup{Argmin}\lk f(t):t\in T\rk := \lk \bar t\in T : f(\bar t) \le f(t) \mbox{ for all }t\in T\rk \, .$$
The nonnegative orthant is denoted by $\R^n_+$. For a vector $\x\in \R^n_+$, the index set
$$
\supp \x = \lk i\in [1\!:\!n] : x_i > 0\rk
$$
is the {\em support}  of $\x$.
Let $\e_i$ be the $i$th column vector of the $n\times n$ identity matrix $\Ib_n$ and $\e=\sum_{i=1}^n \e_i$.
The zero vector and the zero matrix (of appropriate sizes) are denoted by $\o$ and $\Ob$, respectively, and
$\Delta = \{\x\in \R^n_+ : \e\T\x = 1 \}$ stands for the standard simplex.
The vector space of real symmetric $n\times n$ matrices is denoted by $\csn$, and the Frobenius inner product of two
matrices $\lk\Ab,\Bb\rk \subset \csn$ by $\langle \Ab,\Bb\rangle := \tr(\Ab  \Bb)$.
For an $n\times p$ matrix $\Vb=[\v_1, \ldots ,\v_p]$, the relation $\Mb =\Vb\Vb\T$ is equivalent to $\Mb =\sum_{i=1}^p \v_i\v_i\T$. We will refer to this sum as a ``cp decomposition" of $\Mb$, if $\Vb$ has no negative entries.
Given a square matrix $\Sb$, we will, by slight abuse of language, use the phrase ``zero(es) of $\Sb $" as an abbreviation of ``zero(es) of the
quadratic form $\qf\x \Sb $ over $\x\in \Delta$"; this terminology differs slightly from that in~\cite{Hild14a}.

By $\cs n$ we denote the cone of completely positive matrices,
$$\cs n= \conv\{\x\x\T :  \x\in \R^n_+ \}\, .$$
Both, $\cs n$ and its dual, the cone of copositive matrices
$$\cc n = \lk \Sb\in \csn : \x\T \Sb\x \ge0 \mbox{ for all }\x\in \R^n_+\rk \, ,$$
are pointed closed convex cones with nonempty interior. The copositive cone $\cc n$ and, in particular, its extremal rays, are important as any matrix on the boundary $\partial \cs n$ of $\cs n$ is orthogonal to an extremal ray of $\cc n$. So, studies of the extremal rays of $\cc n$ like in~\cite{Dick12c,Hild11,Hild14a} 
lead to conclusions on all matrices on $\partial \cs n$, which allow for inference on \emph{upper} bounds on $p_n$. This was an essential ingredient of the arguments in~\cite{Shak13b,Shak13a}. Here we employ a somewhat reverse approach: we start from (appropriate) matrices $\Sb \in \partial \cc n$ and construct
$\Mb\in \partial \cs n$ where we can calculate the cp-rank $\cpr \Mb$, improving upon \emph{lower} bounds on $p_n$. Eventually, this will lead to examples refuting the DJL conjecture.

\section{Iterative reduction of the cp-rank}

Consider a copositive matrix $\Sb \in \partial \cc n$ and assume that $\lk \q_1,\ldots ,\q_m\rk$ are all the zeroes of $\Sb $. Since $\Sb \in \partial \cc n$,  there is a matrix $\Mb \in \cs n\setminus \lk \Ob \rk$ such that $\du \Mb \Sb =0$, e.g., any matrix of the form
$$\Mb =\sum_{i=1}^m y_i \q_i\q_i\T$$
for some $\y\in \R^m_+\setminus \lk \o\rk$. The next result shows the converse of this statement, so that the set of possible cp decompositions of matrices orthogonal to $\Sb $ is quite restricted:

\begin{lem}\label{orthdesc}
Let  $Q=\lk \x\in \Delta : \qf\x\Sb = 0\rk$ be all the zeroes of  $\Sb \in \partial \cc n$. Then any matrix $\Mb \in \cs n$ orthogonal to $\Sb $ must be of the form
\beq{od}\Mb =\sum_{j=1}^m y_j \q_j\q_j\T\quad \mbox{with }\lk \q_1,\ldots , \q_m\rk \subseteq Q\eeq
for some $\y\in \R^m_+$.
\end{lem}
\bep
Let $\Mb $ have the cp decomposition $\Mb =\sum\limits_{i=1}^m \v_i\v_i\T$ with $\v_i\in \R^n_+\setminus\lk \o\rk$ for all $i\in [1\!:\! m]$. Then  $ \Mb \perp \Sb  $ implies
$$0=\du \Mb \Sb  = \sum_{i=1}^m \qf{\v_i}\Sb \, ,$$
and as $\Sb $ is copositive, every term in above sum must be zero. So all $\q_i := \frac 1{\e\T\v_i }\, \v_i\in  Q$, and the result~\reff{od} follows with $y_i:= (\e\T\v_i)^2$. \ep

Although we have restricted the possible cp decompositions by above observation, there still could be infinitely many, but they can be obtained in a linear way. To be more precise, suppose that $Q=\lk \q_1,\ldots, \q_m\rk$, fix any $\y\in \R^m_+$  such that~\reff{od} holds,  and consider
\beq{defpy}
X_\y := \lk \x\in \R^m_+: \sum_{i=1}^m x_i\q_i\q_i\T = \sum_{j=1}^m y_j\q_j\q_j\T \rk\, .
\eeq
A particular case is obtained if $X_\y = \lk \y\rk$, because then $\cpr \Mb  = | \supp \y |$ is immediate from~Lemma~\ref{orthdesc}. However, this may not always be the case, but
the next theorem will show how to fix some variables $x_k$ of points $\x\in X_\y$ to $y_k$, with some consequences for the construction of matrices of high cp-rank. To apply that theorem in more general situations, we need some further notation. First define for $Q\subseteq\Delta$ the set $\overline Q:=\lk \q\q\T : \q\in Q\rk \subset \csn$ and by
$\cone\! \overline Q := \R_+\conv \overline Q$ the {\em convex conic hull} of $\overline Q$; moreover, for finite $P\subset \Delta$, we denote by
$$\raisebox{3.5pt}{${\circ\atop\cone}$} \!\! \overline P := \lk \sum_{\f\in P} y_\f \f{}\f\T : y_\f>0 \mbox{ for all }\f\in P\rk\, .$$
Finally, we abbreviate the set of completely positive matrices whose cp decompositions can only use multiples of vectors from $Q$ by
$$\mathcal{E}(Q) := \lk \Mb\in \cone\! \overline Q : \mbox{if }\,\Mb\in \raisebox{3.5pt}{${\circ\atop\cone}$} \!\! \overline P \textup{ for finite }P\subseteq\Delta\, , \textup{ then } P\subseteq Q\rk \, .$$
So Lemma~\ref{orthdesc} would read: If $Q$ is the set of all zeroes of $\Sb\in\cc{n}$, then any $\Mb\in\cs{n}$ with $\langle\Mb,\Sb\rangle=0$ satisfies $\Mb\in\mathcal{E}(Q)$.

\begin{thm}\label{varfix}
For a finite subset $Q\subset\Delta$ consider $\Mb\!=\!\sum\limits_{\f\in Q} y_\f\f{}\f\T$ with $\y\!\in\!\mathbb{R}^{|Q|}_+$ and assume $\Mb\in\mathcal{E}(Q)$.
Suppose that there is $\q\in Q$ such that for two different indices $r,s$, we have
\beq{plus}
\lk r,s\rk \subseteq\supp {\q} \quad\mbox{but}\quad \lk r,s\rk \not\subseteq\supp {\q'}\mbox{ for all }\q'\in Q\setminus\{\q\} \, .
\eeq
Then
\begin{enumerate}
  \item [{\rm(a)}] $x_\q  = \frac {\e_r\T\Mb\e_s}{(\e_r\T\q)(\e_s\T\q)}= y_\q$ holds for all $\x\in X_\y$,
  \item [{\rm(b)}] $\widehat\Mb{}: = \Mb  - y_\q\q\q\T\in\mathcal{E}\left(Q\setminus\{\q\}\right)$,
  \item [{\rm(c)}] $\cpr \Mb  = \sgn(y_\q) + \cpr \widehat\Mb$.
\end{enumerate}
\end{thm}

\bep
Condition~\reff{plus} implies $(\e_r\T\q)(\e_s\T\q)>0$ and further that
$$x_k (\e_r\T\q)(\e_s\T\q)=\e_r\T\Mb\e_s\quad \mbox{for all }\x\in X_\y\,.$$
Hence $x_k= \frac {\e_r\T\Mb\e_s}{(\e_r\T\q)(\e_s\T\q)}$ is fixed, which proves (a).
Now define $\widehat y_\q=0$ and $\widehat y_{\q'}=y_{\q'}$ for $\q'\in Q\setminus\{\q\}$, and observe $\widehat \Mb=\sum_{\f\in Q}\widehat y_\f\f{}\f\T\in\mathcal{E}(Q)$. Assertion (a), applied to $\widehat\y$, tells us $\widehat x_\q=\widehat y_\q=0$ for all $\widehat\x\in X_{\widehat y}$, therefore (b) holds.
By (a), any minimal cp decomposition of $\Mb$ is of the form $\Mb=y_\q\q\q\T+\widehat\Mb$, which implies (c).
\ep


If the hypotheses of Theorem~\ref{varfix} including condition~\reff{plus}
are satisfied for $\Mb,Q,\q$, then by (b) of that theorem we know that  $\widehat\Mb\in \mathcal{E}(\widehat Q)$ for  $\widehat Q:=Q\setminus\{\q\}$.
Now, if we want to apply Theorem~\ref{varfix} iteratively, then we may replace $Q$ with $\widehat Q$ so that condition~\reff{plus} may be satisfied more easily for some $\widehat \q\in \widehat Q$.

So if we arrange the supports of (many) $\q$'s such that condition~\reff{plus}, or a similar one, continues to hold during the iterations, we can construct $\Mb$ with high $\cpr \Mb$, as long as $y_\q>0$ continues to hold, too. This will be done in the next section.



\section{Zeroes of cyclically symmetric matrices}

We will employ symmetry transformations of the coordinates given by cyclic permutation, denoting by $a\oplus b$ and $a\ominus b$ the result of addition and subtraction modulo $n$. To keep in line with previous and standard notation, we consider the remainders $[1\!:\! n]$ instead of $[0\!:\! n-1]$, e.g. $1\oplus (n-1) = n$. To be more precise, let
$\Pb_i$ be the square $n\times n$ permutation matrix which effects $\Pb_i\x = [x_{i\oplus j}]_{j\in [1:n]}$ for all $\x\in \R^n$ (for example, if $n=3$ then
$\Pb_2\x = [x_3,x_1,x_2]\T$). Obviously $\Pb_i=(\Pb_1)^i$ for all  integers $i$ (recall $\Pb^{-3}$ is the inverse matrix of $\Pb\Pb\Pb$), $\Pb_i\T = \Pb_{n-i}= \Pb_i^{-1}$ and $\Pb_n= \Ib_n$. A {\em circulant matrix} $\Sb=\Cb(\a)$ based on a vector $\a\in \R^n$ (as its last column rather than the first)
is given by
$$\Sb=[\Pb_ {n-1}\a,\Pb_ {n-2}\a, \ldots, \Pb_ 1\a, \a]\, .$$
If $\Sb=\Cb(\a)\in \csn$, i.e., if $\Cb(\a)$ is symmetric, it is called {\em cyclically symmetric}.

\begin{lem}\label{zyksym}
Any circulant matrix $\Sb = \Cb(\a)$ satisfies $\Pb_i\T\Sb\Pb_i=\Sb$ for all $i\in [1\!:\! n]$. Furthermore, if
\beq{zysy}a_i=a_{n-i}\quad\mbox{for all }i\in[1\!:\! n-1]\, ,\eeq
then $\Sb=\Cb(\a)$ is cyclically symmetric.
\end{lem}
\bep The first relation is evident. To show the remaining assertion, assume~\reff{zysy} and
let $\e_j\T\Sb\e_i=\e_j\T \Pb_ {n-i}\a   = a_k$ with $k\oplus i = j$ while
$\e_i\T\Sb\e_j= a_{ \ell }$ with $ \ell \oplus j=i$. Thus $i\oplus j = k\oplus \ell \oplus i\oplus j$ and $\lk k,\ell\rk\subseteq[1\!:\! n]$, so we get $k+\ell\in \lk n,2n\rk$ and therefore $a_k=a_\ell$. Hence $\Cb(\a)\in \csn$. \ep

Copositive cyclically symmetric matrices $\Sb=\Cb(\a) \in \partial\cc n$ can have many zeroes (which then are global minimizers
of the quadratic form $\qf\x \Sb$ over $\Delta$; for local minimizers this has already been observed earlier, see~\cite{Broo93a} and references therein). To facilitate the argument, let us denote by $\Rb\in \csn$ the {\em reflection matrix} which transforms every $\x\in \R^n$ into its {\em mirror image} $\Rb\x:=[x_{n+1-i}]_{i\in [1:n]}$. Note that $\Rb\T=\Rb\in \csn$.

In the sequel, it will be convenient to denote, for any $\q\in\mathbb{R}^n_+$, the set $D_\q$ of differences and the set $U_\q$ of unique differences of the elements of $I(\q)$:\\
\hspace*{0.5cm}$D_\q\! :=\!\{d\!\in\![1\!:\!n\!-\!1]:\, d\!=\!r\!\ominus\! s \textup{ has at least one solution with \!}\lk r\!,\!s\rk\! \subseteq\!\supp \q \}\,, $\\
\hspace*{0.5cm}$U_\q\! :=\!\{d\!\in\![1\!:\!n\!-\!1]:\, d\!=\!r\!\ominus\! s \textup{ has exactly one solution with \!}\lk r\!,\!s\rk\! \subseteq\supp \q \}\,. $


\begin{lem}\label{manyzero}
Let $\Sb=\Cb(\a)\in \csn$ be a cyclically symmetric matrix.
\begin{enumerate}
  \item[{\rm(a)}]
We have $\Rb \Sb\Rb = \Sb$. Further, fixing $\q\in\mathbb{R}^n_+$,
for any shift $\q'=\Pb_i\q$, and for its mirror image $\q''=\Rb\q$, we have
\beq{invzero}
(\q')\T \Sb\q'  = (\q'')\T \Sb\q''=\qf\q \Sb\, .
\eeq
\item[\rm{(b)}]
For any zero $\q$ of  $\Sb$ there are actually up to $2n$ zeroes: the shifts $\Pb_i\q$ for $i\in[1:n]$ and their mirror images, if they are all different.
 \item[\rm{(c)}] The supports of zeroes are shifted cyclically, $\supp {\Pb_i\q} =  \lk j\ominus i : j\in \supp \q\rk$.
 However, the relative  differences within the support of course remain: if $\lk r,s\rk \subseteq \supp \q$, then
$r\ominus s = r'\ominus s'$ if $r'=r\oplus i$ and $s'=s\oplus i$.
 \item[\rm{(d)}] For any $\q\in\mathbb{R}^n_+$ the sets $D_\q$ and $U_\q$
 are invariant under shifts and reflection: with $\q',\,\q''$ as in (a), we have $D_\q\!=\!D_{\q'}\!=\!D_{\q''}$ and $U_\q\!=\!U_{\q'}\!=\!U_{\q''}$.\\
 Moreover $d=\frac n2\not\in U_\q$ for even $n$, since then $r\!\ominus\! s\!=\!d$ implies $s\!\ominus\! r\!=\!d$.
\end{enumerate}
\end{lem}
\bep The relation $ \Rb \Sb\Rb = \Sb$ can be checked in a straightforward manner while the equations in~\reff{invzero} follow from
$$(\q')\T \Sb\q'  =\q\T\Pb_i\T\Sb\Pb_i\q  =\qf\q \Sb$$
and from
$$(\q'')\T \Sb\q''   =\q\T\Rb\T\Sb\Rb\q  =  \q\T\Rb\Sb\Rb\q =  \qf\q \Sb\, .$$
The assertions about the supports are evident.
\ep

\begin{thm}\label{reldi}
For a finite subset $Q\subset\Delta$ consider $\Mb\!=\!\sum\limits_{\f\in Q} y_\f\f{}\f\T$ with $\y\!\in\!\mathbb{R}^{|Q|}_+$ and assume $\Mb\in\mathcal{E}(Q)$.
Fix $\q\in Q$, define $Q_1:=\{\Pb_i\q:i\in[1\!:\! n]\}$ and define $Q_2:=Q\setminus Q_1$. Assume $Q_1\subseteq Q$, and that there is $d\in[1\!:\! n-1]$ such that $d\in U_{\q}\setminus\bigcup\limits_{\q'\in Q_2}D_{\q'}$. Then
\begin{enumerate}
  \item [{\rm(a)}] $x_{\f}  = y_{\f}$ holds for all $\x\in X_\y$ and $\f\in Q_1$,
  \item [{\rm(b)}] $\widehat\Mb{}: = \Mb  - \sum\limits_{\f\in Q_1}y_\f\f{}\f\T
     \in\mathcal{E} \left(Q_2\right)$,
  \item [{\rm(c)}] $\cpr \Mb  = \sum\limits_{\f\in Q_1}\sgn(y_\f) + \cpr \widehat\Mb$.
\end{enumerate}
\end{thm}

\bep
Let $\{r,s\}\subseteq I(\q)$ satisfy $d=r\ominus s$.
By the assumptions it is clear that $\lk r,s\rk\subseteq\supp{\q'}$ can never hold for any $\q' \in Q_2 $. So consider instead $ \f =\Pb_i\q$ for $i\in [1\!:\! n-1]$. We argue by contradiction: if $\lk r,s\rk\subseteq\supp{ \f }$, then  $\lk r\oplus i,s\oplus i\rk\subseteq\supp {\q} $ but differs from the pair $\lk r,s\rk$
(note that $r=s\oplus i$ and simultaneously $s=r\oplus i$ is impossible since $d\neq\frac n2$).
Obviously the difference would be the same, namely $d$, which by assumption is absurd. So condition~\reff{plus} holds for $Q$ and $\q$.
Since $U_{\f}=U_{\q}$ for all $\f\in Q_1$, by Lemma~\ref{manyzero}~(d), we similarly obtain that condition~\reff{plus} holds for $Q$ and $\f\in Q_1\setminus\{\q\}$. Finally we obtain (a), (b) and (c) by iterating the reduction step of Theorem~\ref{varfix} in total $|Q_1|$ times.
\ep


\begin{cly}
Let all hypotheses of Theorem~\ref{reldi}  be satisfied with $Q_2=\emptyset$. Let $\Mb :=\sum\limits_{\f\in Q}y_\f \,\f\,\!\f\T$.
If $y_\f> 0$ for all $\f\in Q$, then the minimal cp decomposition of $\Mb$ is unique and $\cpr \Mb=| Q|$.
\end{cly}

The next two results deal with instances where there is more than one minimal cp decomposition of a similarly constructed matrix:

\begin{lem}\label{rankdef}
Consider $\q\in\mathbb{R}^n_+$ such that $Q:=\{\Pb_i\q:i\in[1\!:\! n]\}$ satisfies $|Q|=n$ and $\Rb\q\notin Q$.
Suppose there are $d_1,d_2\in\ U_\q$  with $d_1=r\ominus s$ and $d_2=\rho\ominus \sigma$, such that $\rho+\sigma-r-s$ and $n$ are coprime.  We consider the following subset of $\csn$:
$$\mathcal{F}:=\lk \f\,\!\f\T:\f\in Q\rk \cup\{\Rb\f (\Rb\f)\T:\f\in Q\}\, .$$
Then every $(2n-1)$-element subset of $\mathcal{F}$ is linearly independent, moreover $\mathcal{F}$ itself (as a subset of the vector space $\csn$) has rank $2n-1$.
\end{lem}

\bep
We first observe that our assumptions on $Q$ imply $|\mathcal{F}|=2n$.
Moreover, $ U_{\Rb\q}=U_\q$.
We now claim that
\beq{ffzysy}\sum_{\f\in Q }\Rb\f (\Rb\f)\T=\Rb\left(\sum_{\f\in Q }\f\,\!\f\T\right) \Rb
=\sum_{\f\in Q }\f\,\! \f\T\, .\eeq
The last equality can be established in the following way. Note that $\Ab:=\sum_{\f\in Q }\f\,\! \f\T$ can be rewritten as $\Cb (\a )$ with $a_i =\q\T \Pb_i \q$, because
$$\e_r\T \left [ \sum_{i=1}^n (\Pb_i\q)(\Pb_i\q)\T \right]\e_s = \sum_{i=1}^n q_{i\oplus r} q_{i\oplus s} = \sum_{i=1}^n q_i q_{i \oplus r \ominus s}= \q \T \Pb_{r \ominus s} \q $$
depends on $(r,s)$ only via $r\ominus s$. Symmetry of $\Ab=\Cb(\a)$ follows from Lemma~\ref{zyksym} because condition~\reff{zysy} is satisfied due to
$$ a_i = \q\T \Pb_i \q=\q \T \Pb_{n-i}^{-1} \Pb_{n-i} \Pb_i \q= \q \T \Pb_{n-i}\T \q = a_{n-i}\quad \mbox{for all }i\in [1\!:\! n-1]\, .$$
Equality~\reff{ffzysy} is now  established by Lemma~\ref{manyzero}(a).
Hence the rank of $\mathcal{F}$ can be at most $2n-1$.
Let $\q_i:=\Pb_i\q$ for $i\in[1\!:\! n]$. Then $\{r\ominus i,s\ominus i\}\subseteq\supp{\q_i}$. Further define $\q'_i:=\Rb\q_i=\Rb\Pb_i\q=\Pb_{n-i}\Rb\q$ and note that $\e_a\T\q_b= q_{a\oplus b}$ as well as $\e_a\T \q'_b= q_{1 \oplus b\ominus a}$ for all $a,b\in [1\!:\! n]$.
Next consider the equation
\beq{linind}
\sum_{i=1}^n x_i\q_i\q_i\T +\sum_{i=1}^n x'_i\q'_i{\q'_i}\T=\Ob.
\eeq
Multiplying with $\e_{r\ominus j}\T$ from the left and with $\e_{s\ominus j}$ from the right, we obtain
$$\sum_{i=1}^n x_i q_{r\ominus j \oplus i} q_{s \ominus j \oplus i} +\sum_{i=1}^n x'_i q_{1\oplus i \ominus \left ( r \ominus j \right )} q_{1 \oplus i \ominus \left ( s \ominus j \right )}=0.$$
By the assumptions on $U_{\q}$ we see that the only terms contributing to the sum are achieved by
choosing $i=j$ in the first term and $i=  r \oplus s\ominus j \ominus 1$ in the second term. This results in
\addtocounter{equation}{1}
\beq{linsyspart1}x_j q_r q_s + x'_{r \oplus s\ominus j \ominus 1} q_s q_r =0 \quad \textup{ for all }j\in[1:n]\, .
\tag*{(\theequation{a})}\eeq
Similarly multiply with $\e_{\rho\ominus j}\T$ from the left and with $\e_{\sigma\ominus j}$ from the right, yielding
\beq{linsyspart2}x_j q_\rho q_\sigma + x'_{\rho \oplus \sigma\ominus j \ominus 1} q_\sigma q_\rho =0 \quad \textup{ for all }j\in[1\!:\! n]\, .\tag*{(\theequation{b})}\eeq
From these equations we conclude that $x_j'=x'_{j\oplus \rho\oplus \sigma\ominus r\ominus s} =x'_{j\oplus (\rho+\sigma-r-s)}$ for all $j\in[1\!:\! n]$.  Fixing $x_1'=\xi$, and employing coprimality of $\rho+\sigma-r-s$ and $n$, we see that our system $(\theequation)$ of $2n$ equations has the unique solution $x_i=-x_i'=-\xi$ for $i\in[1\!:\! n]$. So there is a one parameter family of solutions parameterized by $\xi$, showing that if any of the coefficients in \eqref{linind} is zero, all others also must be zero, so indeed every $(2n-1)$-element subset of $\mathcal{F}$ has  to be linearly independent, as asserted.
\ep

\begin{thm}\label{rankdefcly}
Let $\q$ satisfy the hypotheses of Lemma~\ref{rankdef} and define the (therefore disjoint) sets $Q:=\{\Pb_i\q:i\in[1\!:\! n]\}$ and $Q':=\{\Rb\q:\q\in Q\}$.
Consider the matrix $\Mb = \!\! \!\! \sum\limits_{\f\in Q\cup Q'}\!\! y_\f\,\f\,\!\f\T$ with
$\y\in\mathbb{R}_+^{|Q\cup Q'|}$, and
assume that $\Mb\in\mathcal{E}(Q\cup Q')$ holds. Then we have:
\begin{enumerate}
  \item[{\rm(a)}]  If all $y_\f>0$ and if $|\textup{Argmin} \lk y_\f:\f\in Q\rk|=|\textup{Argmin} \lk y_\f:\f\in Q'\rk|=1$,
then there are exactly two different minimal cp decompositions of $\Mb$ and $\cpr \Mb =2| Q |-1$.
  \item[{\rm(b)}] If $y_\f=0$ for at least one $\f\in Q$ and at least one $\f\in Q'$, then the minimal cp decomposition of $\Mb$ is unique and
  $\cpr \Mb =| \supp \y |$.
  \end{enumerate}
\end{thm}
\bep
Define $u_\f:=1$ for all $\f\in Q$ and $u_\f:=-1$ for all $\f\in Q'$. Then, by Lemma~\ref{rankdef} and equation~\reff{ffzysy}, the solutions $\x$ of the equation
$\Mb = \sum\limits_{\f\in Q\cup Q'}x_\f\,\f\,\!\f\T$ are given by $\x=\y+\xi\u$.
In case (a), the solutions $\x\ge\o$ additionally require $\xi\in[-\min\{y_\f:\f\in Q\},\min\{y_\f:\f\in Q'\}]$, with $|I(\x)|=2|Q|-1$ (resp.~$|I(\x)|=2|Q|$) for $\xi$ on the boundary (resp.~in the interior) of that interval. In case (b), the condition $\x\ge\o$ is violated for any $\xi\ne0$, so $\x=\y$ is unique.
\ep

\section{Counterexamples to the Drew-Johnson-Loewy conjecture}

For the examples to follow, we selected matrices $\Sb$ with integer entries, where we could determine all minimizers of the quadratic form $\qf\x\Sb$ by exact arithmetic, solving the first-order conditions and checking the values for nonnegativity with the help of~\reff{invzero}, cf. also~\cite{Bomz92a,Bomz02a}. To be more precise, we first checked (by exact arithmetic to avoid any numerical errors) for all possible supports $I\subseteq [1:n]$ with $I\neq\lm$, whether there could be a {\em local} solution $\q$ to the optimization problem $\min\limits_{\q\in \Delta} \qf\q\Sb$ with $I(\q)=I$. To this end, ignoring the variables $q_i=0$ for $i\in [1\!:\! n]\setminus I$, we see there is only one locally binding constraint, namely $\e\T\q=1$. So, denoting the multiplier of this constraint by $2m$, we arrive at the first-order conditions
\beq{helpsys}\left. \bea{rcl}  \Sb_I\x &= &m\e  \\[0.2em]
                                \e\T\x &= &1  \ea\right \}\, , \eeq
where $\Sb_I$ denotes the principal submatrix of $\Sb$ on $I\times I$. Since all constraints of the optimization problem $\min\limits_{\q\in \Delta} \qf\q\Sb$ are linear, any local minimizer of $\qf\q\Sb$ over $\Delta$ must solve~\reff{helpsys} for some $I$, putting $\x=[q_i]_{i\in I}$ and $m = m\x\T\e = \qf\x{\Sb_I} = \qf\q\Sb$.
Now suppose $\e\T\x=\e\T\y=1$ and $\Sb_I\x=m\e$ and $\Sb_I\y=t\e$. Then
\beq{unval} t= ( t\e)\T\x = \y\T\Sb_I\T \x = \y\T \Sb_I\x = \y\T (m\e) = m \, ,\eeq
so that the value $m = \qf\x{\Sb_I} =: m_I$ at any solution $(m,\x)\in \R\times \R^{|I|}$ to~\reff{helpsys} is uniquely determined by $I$.
We solved~\reff{helpsys} by exact arithmetic for all $I$. If there is a unique solution $(m_I,\x_I)$, we discarded $I$ where $\x_I \notin \R^{|I|}_+$. For the remaining $I$, we confirmed that $m_I\ge0$ if~\reff{helpsys} has a solution at all. This established copositivity of $\Sb$.
The next step is to determine all zeroes of $\Sb$, i.e., all solutions $(0,\x)$ to~\reff{helpsys} with $\x\in \R^{|I|}_+$. While there could be multiple solutions to~\reff{helpsys} for $m_I>0$, this is ruled out in case $m_I=0$ for the matrices $\Sb$ considered below.
Indeed, consider again two solutions $(0,\x)$ and $(0,\y)$ to~\reff{helpsys}. Then $\Sb_I(\x-\y) = \o$  and $\x-\y\in \ker \Sb_I \cap \e ^\perp$, so that the condition $\ker \Sb_I\cap \e^\perp=\lk \o\rk$ rules out multiple solutions to~\reff{helpsys}; this is in fact true for any value of $m_I$, due to~\reff{unval}.
Now \cite[Lemma~1]{Bomz02a} shows that $\ker \Sb_I\cap \e^\perp=\lk \o\rk$ holds if $\e\e\T- \Sb_I$ is nonsingular, which we confirmed (again by exact arithmetic) for all $I$ which admit a solution $(0,\x)$ to~\reff{helpsys} with $\x\in \R^{|I|}_+$.
Note that if $(0,\x)$ solves~\reff{helpsys}, then $(\e\e\T-\Sb_I)\x= \e$ and $\e\T\x=1$ implies $\e\T (\e\e\T-\Sb_I)^{-1} \e=1$. So we considered the unique solution
$\x_I := (\e\e\T-\Sb_I)^{-1} \e\in \R^{|I|}_+$. Finally, we filled the entries with indices in $[1\!:\! n]\setminus I$ by zeros to get a vector which we call $\q_I\in \R^n$ and collected these as the set of all zeroes of $\Sb$. In this way the assumptions of the previous sections were ensured. As a final remark, note that \cite[Lemma~1]{Bomz02a} says that $\ker \Sb_I\cap \e^\perp\neq\lk \o\rk$ holds if and only if {\em both} $\Sb_I$ {\em and} $\Sb_I-\e\e\T$ are singular; however, as noted by the Senior Editor, in case $m_I=0$  the principal submatrices $\Sb_I$ necessarily have to be singular as they must be positive-semidefinite (see, e.g.~\cite[Lemma~2.4]{Dick12c}), so the above argument involving $\Sb_I-\e\e\T$ is essential.

\textbf{Example 1 $(p_7\ge14)$:} Let $\Sb=\Cb([-153,127,-27,-27,127,-153,162]\T)$. Then the set of zeroes of $\Sb$ in $\Delta \subset\mathbb{R}^7$ consists of 14 vectors: 
$\q_i=\Pb_i\u, i\in[1\!:\! 7]$, where $\u=\frac17[3,3,0,0,1,0,0]\T$,
and $\q_i=\Pb_i\v, i\in[8\!:\! 14]$, where $\v=\frac1{35}[9,17,9,0,0,0,0]\T$.
Let
$$\Mb:=\sum_{i=1}^{14}\q_i\q_i\T={\textstyle\frac1{1225}}\,\Cb([531,81,150,150,81,531,926]\T).$$
The difference sets of $\supp \u =\{1,2,5\}$ and $\supp \v =\{1,2,3\}$ are
$$D_\u=\{1,3,4,6\},\quad U_\u=\{1,6\},\quad D_\v=\{1,2,5,6\},\quad U_\v=\{2,5\}.$$
We note that $d=2\in U_\v\setminus D_\u$, so we may apply Theorem~\ref{reldi} with $Q=\{\q_1,\ldots,\q_{14}\}$, $\q=\v$ and
$Q_1=\{\q_8,\ldots,\q_{14}\}$, to conclude that in any cp decomposition
$\Mb=\sum\limits_{i=1}^{14}x_i\q_i\q_i\T$ we must have $x_k=1$ for $k\in[8\!:\! 14]$.
Moreover Theorem~\ref{reldi} states that
$\widehat\Mb:=\Mb-\sum\limits_{i=8}^{14}\q_i\q_i\T=\sum\limits_{i=1}^{7}x_i\q_i\q_i\T$ satisfies $\widehat\Mb\in\mathcal{E}(\widehat Q)$, where
$\widehat Q=\{\q_1,\ldots,\q_{7}\}$.
Therefore $d=1\in U_{\u}$ allows us to invoke Theorem~\ref{reldi} with
$\Mb=\widehat\Mb$, $Q=\widehat Q$ and $\q=\u$.
 We conclude  that $x_k=1$ also for $k\in[1\!:\! 7]$, that $\Mb$ has a unique minimal cp decomposition, and that $\cpr \Mb=14$. Another matrix of this sort, having small integer entries, is
$$ \widetilde M_7:={\tfrac23} 7^2\sum_{i=1}^{7}\q_i\q_i\T+{\tfrac13} 35^2\sum_{i=8}^{14}\q_i\q_i\T
 =
 \small{\left [
  \begin{array}{ccccccc}
  163&108&27&4&4&27&108\\108&163&108&27&4&4&27\\27&108&163&108&27&4&4\\4&27&108&163&108&27&4\\4& 4&27&108&163&108&27\\27&4&4&27&108&163&108\\108&27&4&4&27&108&163
  \end{array}
\right]}
 \, .
$$
Note that both, above matrix and $\Mb$, have no zero entries and full rank.

\vspace{.25cm}%

\textbf{Example 2 $(p_9\ge26)$:}
Let
$$\Sb=\Cb([-1056, 959, -484, 231, 231, -484, 959, -1056, 1089]\T).$$
Then the set of zeroes of $\Sb$ in $\Delta\in\mathbb{R}^9$ consists of 27 vectors:  indeed, let
$$\left . \bea{ccc}
\!\!\u\!\! &=\!\! &\frac1{26}[11,12,0,0,3,0,0,0,0]\T \\[0.2em]
\!\!\v\!\! &=\!\! &\frac1{26}[12,11,0,0,0,0,3,0,0]\T\\[0.2em]
\!\!\w \!\!&= \!\!&\frac1{130}[33,64,33,0,0,0,0,0,0]\T\ea\!\!\!\!\rk\mbox{ and define }\; \q_i := \lk \bea{ll}
\!\!\Pb_i\u\, ,\!\! &\mbox{if }i\in[1\!:\! 9]\, ,\\[0.2em]
\!\!\Pb_i\v\, , \!\! &\mbox{if }i\in[10\!:\! 18]\, ,\\[0.2em]
\!\!\Pb_i\w\, ,\!\!   &\mbox{if } i\in[19\!:\! 27]\, .
\ea\right .$$
The set of zeroes of $\Sb$ is $\lk \q_i : i\in [1\!:\! 27]\rk$ and $\Pb_2\v=\Rb\u\notin\lk \Pb_i\u:i\in [1\!:\! 9]\rk$. Put
\begin{align*}
\Mb:=&2\sum_{i=1}^{18}\q_i\q_i\T-\q_{9}\q_{9}\T-\q_{11}\q_{11}\T+\sum_{i=19}^{27}\q_i\q_i\T\\
=&{\tfrac1{16900}}
\scriptsize{
\left [
  \begin{array}{ccccccccc}
30649&14124&1089&3600&2475&3300&3600&1089&17424\\14124&30074&17424&1089&2700&3300& 3300&3600&1089\\1089&17424&33674&17424&1089&3600&3300&3300&3600\\3600&1089&17424& 33674&17424&1089&3600&3300&3300\\2475&2700&1089&17424&33224&17424&1089&2700& 2475\\3300&3300&3600&1089&17424&33674&17424&1089&3600\\3600&3300&3300&3600&1089& 17424&33674&17424&1089\\1089&3600&3300&3300&2700&1089&17424&30074&14124\\17424&1089& 3600&3300&2475&3600&1089&14124&30649\\
   \end{array}
\right]}.
\end{align*}
\normalsize
The difference sets of $\supp \u =\{1,2,5\}$, $\supp\v =\{1,2,7\}$ and $\supp\w =\{1,2,3\}$ are
$$D_\u=U_\u=D_\v=U_\v=\{1,3,4,5,6,8\},\quad D_\w=\{1,2,7,8\},\quad U_\w=\{2,7\}.$$
We note that $d=2\in U_\w\setminus (D_\u\cup D_\v)$, so we may apply Theorem~\ref{reldi} to conclude that in any cp decomposition
$\Mb=\sum\limits_{i=1}^{27}x_i\q_i\q_i\T$ we must have $x_i=1$ for $i\in[19\!:\! 27]$. Next, consider the matrix
$\widehat\Mb:=\Mb-\sum\limits_{i=19}^{27}\q_i\q_i\T$, which satisfies $\widehat\Mb\in\mathcal{E}(\widehat Q)$ by Theorem~\ref{reldi}, where $\widehat Q=\{\q_1,\ldots,\q_{18}\}$.
Since the differences $d_1=1=2-1$ and $d_2=3=5-2$ appear only once in $U_{\u}$,
 and $5+2-2-1=4$ and $9$ are coprime allows to invoke Lemma~\ref{rankdef} and Theorem~\ref{rankdefcly} with $\Mb=\widehat \Mb$, $Q\cup Q'=\widehat Q$ and $\q=\u$.
 We conclude that there are exactly two vectors $\x\in\mathbb{R}^{18}_+$ of support of size 17, (and no such vectors of smaller support,) that give rise to minimal cp decompositions of $\widehat \Mb$, and that $\cpr \Mb=26$.
 Another matrix of this sort, having small integer entries, is
\begin{align*}
\widetilde \Mb_9:=& {\tfrac56 \, 26^2}\left(\sum_{i=1}^{18}\q_i\q_i\T-{\tfrac35}\,(\q_{7}\q_{7}\T+\q_{13}\q_{13}\T)\right)
 +{\tfrac13}\, 130^2\sum_{i=19}^{27}\q_i\q_i\T\\
=&\small{
\left [
  \begin{array}{ccccccccc}
      2548&1628&363&60&55&55&60&363&1628\\1628&2548&1628&363&60&55&55&60&363\\363&1628&2483&1562&363&42&22&55&60\\60&363&1562&2476& 1628&363&42&55&55\\55&60&363&1628&2548&1628&363&60&55\\55&55&42&363&1628&2476&1562&363&60\\60&55&22&42&363&1562&2483&1628& 363\\363&60&55&55&60&363&1628&2548&1628\\1628&363&60&55&55&60&363&1628&2548\\
   \end{array}
\right]}.
\end{align*}

Note that neither of these matrices of cp-rank 26 are cyclically symmetric, they have no zero entries and full rank.

\vspace{.25cm}%

\textbf{Example 3 $(p_8\ge18)$:}
Continuing Example~2, we observe that the upper left $8\times8$-submatrix $\Sb_8$ of $\Sb$ has 18 zeroes. These are obtained by taking the first 8 coordinates of those zeroes $\q$ of $\Sb$ satisfying $\e_9\T\q=0$. Indeed, if $\z\T \Sb_8\z = 0$, then also $[\z\T | 0]\Sb [\z\T | 0]\T = 0$, so that the zero $ [\z\T | 0]\T $ must appear in the list of Example~2. Define the set $S_8:=\{\q\in\{\q_1,\ldots,\q_{27}\}:\e_9\T\q=0\}$.
Then the matrix $\Mb:=\sum\limits_{\q\in S_8}\q\q\T$ satisfies $\cpr \Mb=18$, by Theorem~\ref{reldi},
Lemma~\ref{rankdef} and Theorem~\ref{rankdefcly}. Moreover all entries in the last row and the last column of $\Mb$ are zero, therefore also $\Mb_8$, the upper left $8\times8$-submatrix of $\Mb$, has $\cpr \Mb_8=18$. Again, by adjusting weights, we came up with a matrix with small integer entries:
$$\widetilde \Mb_8:=
\small{\left [
  \begin{array}{ccccccccc}
541&880&363&24&55&11&24&0\\880&2007&1496&363&48&22&22&24\\363&1496&2223&1452&363&24&22&11\\24&363&1452&2325&1584&363&48&55\\55&48& 363&1584&2325&1452&363&24\\11&22&24&363&1452&2223&1496&363\\24&22&22&48&363&1496&2007&880\\0&24&11&55&24&363&880&541\\
   \end{array}
\right]}.$$
Note that $\widetilde \Mb_8$ is, again, not cyclically symmetric, and that it has full rank.

\vspace{.25cm}%

\textbf{Example 4 $(p_{10}\ge27)$:} Continuing Example~2, let $\Mb\in \cs {10}$ be the matrix obtained from $\widetilde \Mb _9$ by appending a zero column $\o\in \R^9$ and completing this to a symmetric $10\times 10$ matrix by adding one row $\e_{10}\T$ as the last one. Then, by~\cite[Prop.2.2]{Shak13a}, we get
$$\cpr \Mb = \cpr \widetilde \Mb _9 + 1 = 27\, .$$


\textbf{Example 5 $(p_{11}\ge32)$:} Consider
$$\Sb=\Cb([32, 18, 4, -24, -31, -31, -24, 4, 18, 32, 32]\T).$$
There are 33 zeroes of $\Sb$; indeed, let
$$\left . \bea{ccc}
\!\!\u\!\! &=\!\! &\frac1{21}[8, 0, 3, 0, 0, 0, 10, 0, 0, 0, 0]\T \\[0.2em]
\!\!\v\!\! &=\!\! &\frac1{21}[10, 0, 0, 0, 3, 0, 8, 0, 0, 0, 0]\T\\[0.2em]
\!\!\w \!\!&= \!\!&\frac1{7}[2, 0, 0, 2, 0, 0, 0, 3, 0, 0, 0]\T\ea\!\!\!\!\rk\mbox{ and define }\; \q_i := \lk \bea{ll}
\!\!\Pb_i\u\, ,\!\! &\mbox{if }i\in[1\!:\! 11]\, ,\\[0.2em]
\!\!\Pb_i\v\, , \!\! &\mbox{if }i\in[12\!:\! 22]\, ,\\[0.2em]
\!\!\Pb_i\w\, ,\!\!   &\mbox{if } i\in[23\!:\! 33]\, ,
\ea\right .$$
then the set of zeroes can be written as  $\lk \q_i : i\in [1\!:\! 33]\rk$. Now put
\begin{align*}
\Mb:=&2\sum_{i=1}^{22}\q_i\q_i\T-\q_{11}\q_{11}\T-\q_{13}\q_{13}\T+\sum_{i=23}^{33}\q_i\q_i\T\\
=&{\tfrac1{441}}
\scriptsize{
\left [
\begin{array}{ccccccccccc}
 781 & 0 & 72 & 36 & 228 & 320 & 240 & 228 & 36 & 96 & 0 \\
 0 & 845 & 0 & 96 & 36 & 228 & 320 & 320 & 228 & 36 & 96 \\
 72 & 0 & 827 & 0 & 72 & 36 & 198 & 320 & 320 & 198 & 36 \\
 36 & 96 & 0 & 845 & 0 & 96 & 36 & 228 & 320 & 320 & 228 \\
 228 & 36 & 72 & 0 & 781 & 0 & 96 & 36 & 228 & 240 & 320 \\
 320 & 228 & 36 & 96 & 0 & 845 & 0 & 96 & 36 & 228 & 320 \\
 240 & 320 & 198 & 36 & 96 & 0 & 745 & 0 & 96 & 36 & 228 \\
 228 & 320 & 320 & 228 & 36 & 96 & 0 & 845 & 0 & 96 & 36 \\
 36 & 228 & 320 & 320 & 228 & 36 & 96 & 0 & 845 & 0 & 96 \\
 96 & 36 & 198 & 320 & 240 & 228 & 36 & 96 & 0 & 745 & 0 \\
 0 & 96 & 36 & 228 & 320 & 320 & 228 & 36 & 96 & 0 & 845 \\
\end{array}
\right]},
\end{align*}
\normalsize
and again $\Mb$ has full rank.
We get $\supp \u =\{1,3,7\}$, $\supp\v =\{1,5,7\}$, $\supp\w =\{1,4,8\}$, and we calculate
$$D_\u=U_\u=D_\v=U_\v=\{2,4,5,6,7,9\},\quad D_\w=\{3,4,7,8\},\quad U_\w=\{3,8\}. $$
Analogously to Example~2 we now show that the cp-rank is $32$. Since $d=3\in U_\w\setminus (D_\u\cup D_\v)$, we must have $x_i=1$ for $i\in[22\!:\! 33]$ by Theorem~\ref{reldi}.
Therefore consider
$\widehat\Mb:=\Mb-\sum\limits_{i=22}^{33}\q_i\q_i\T$.
We can see that the differences $d_1=6=7-1$ and $d_2=4=7-3$ appear in $U_{\u}$,
 and knowing that $7+3-7-1=2$ and $11$ are coprime allows to invoke Lemma~\ref{rankdef} and Theorem~\ref{rankdefcly}.
 Hence there are exactly two vectors $\x\in\mathbb{R}^{22}_+$ of support of size $21$ for $\widehat \Mb$ and this leads to a total of $32$ for $\cpr \Mb$.


\begin{table}[h]
\caption{(Ranges for) maximal cp-rank $p_n$ of cp matrices of order $n$.}
\label{tab:pnrange}
\begin{center}
\begin{tabular}{| c| ccccccc|}
\hline
   $n$                          & 5               & 6           & 7           & 8            & 9            & 10          & 11 \\[0.2em]
\hline  
$d_n$                      & 6               & 9            & 12            & 16           & 20           & 25           & 30  \\[0.2em]
$p_n$                       & 6              & $\le 15$   & $\ge 14$     & $\ge 18$    & $\ge 26$     & $\ge 27$       & $\ge 32$  \\[0.2em]  
$s_n$                      &11              & 17            & 24            & 32          & 41          & 51             & 62    \\[0.2em]
\hline
\end{tabular}
\end{center}
\end{table}
Table~\ref{tab:pnrange} summarizes the known bracket and consequences from above examples. A tighter upper bound $p_6\le 15$
was proved in~\cite[Thm.6.1]{Shak13b}, but up to now no $\Mb\in \cs 6$ with $\cpr \Mb >9=d_6$ is known.

\vskip 0.5cm
\noindent{\bf{Acknowledgment.}}
The authors are indebted to an anonymous referee for valuable suggestions which significantly improved presentation, and to the Senior Editor Raphael Loewy for his diligence in the evaluation process.

\vskip 0.5cm
\noindent{\bf{References}}

\frenchspacing
\small
\addcontentsline{toc}{section}{References}
\begin{thebibliography}{10}

\bibitem{Berm98}
Abraham Berman and Naomi Shaked-Monderer.
\newblock Remarks on completely positive matrices.
\newblock {\em Linear and Multilinear Algebra}, 44:149--163, 1998.

\bibitem{Berm03}
Abraham Berman and Naomi Shaked-Monderer.
\newblock {\em Completely positive matrices}.
\newblock World Scientific Publishing Co. Inc., River Edge, NJ, 2003.

\bibitem{Bomz92a}
Immanuel~M. Bomze.
\newblock Detecting all evolutionarily stable strategies.
\newblock {\em J. Optim. Theory Appl.}, 75(2):313--329, 1992.

\bibitem{Bomz02a}
Immanuel~M. Bomze.
\newblock Regularity versus degeneracy in dynamics, games, and optimization: a
  unified approach to different aspects.
\newblock {\em SIAM Rev.}, 44(3):394--414, 2002.

\bibitem{Bomz11a}
Immanuel~M. Bomze.
\newblock Copositive optimization -- recent developments and applications.
\newblock {\em European J. Oper. Res.}, 216:509--520, 2012.

\bibitem{Bomz11b}
Immanuel~M. Bomze, Werner Schachinger, and Gabriele Uchida.
\newblock Think co(mpletely~)positive~! -- matrix properties, examples and a
  clustered bibliography on copositive optimization.
\newblock {\em J. Global Optim.}, 52:423--445, 2012.

\bibitem{Broo93a}
Mark Broom, Chris Cannings, and Glenn~T. Vickers.
\newblock On the number of local maxima of a constrained quadratic form.
\newblock {\em Proc. R. Soc. Lond. A}, 443:573--584, 1993.

\bibitem{Bure11}
Samuel Burer.
\newblock Copositive programming.
\newblock In Miguel~F. Anjos and Jean~Bernard Lasserre, editors, {\em Handbook
  of Semidefinite, Cone and Polynomial Optimization: Theory, Algorithms,
  Software and Applications}, International Series in Operations Research and
  Management Science, pages 201--218. Springer, New York, 2012.

\bibitem{Dick12c}
Peter J.~C. Dickinson, Mirjam D\"ur, Luuk Gijben, and Roland Hildebrand.
\newblock Irreducible elements of the copositive cone.
\newblock {\em Linear Algebra Appl.}, 439(6):1605�--1626, 2013.

\bibitem{Drew96a}
John~H. Drew and Charles~R. Johnson.
\newblock The no long odd cycle theorem for completely positive matrices.
\newblock In {\em Random discrete structures}, volume~76 of {\em IMA Vol. Math.
  Appl.}, pages 103--115. 1996.

\bibitem{Drew94}
John~H. Drew, Charles~R. Johnson, and Raphael Loewy.
\newblock Completely positive matrices associated with {$M$}-matrices.
\newblock {\em Linear Multilinear Algebra}, 37(4):303--310, 1994.

\bibitem{Duer10}
Mirjam D{\"u}r.
\newblock Copositive programming --- a survey.
\newblock In Moritz Diehl, Francois Glineur, Elias Jarlebring, and Wim
  Michiels, editors, {\em Recent Advances in Optimization and its Applications
  in Engineering}, pages 3--20. Springer, Berlin Heidelberg New York, 2010.

\bibitem{Hild11}
Roland Hildebrand.
\newblock The extremal rays of the $5\times 5$ copositive cone.
\newblock {\em Linear Algebra Appl.}, 437(7):1538--1547, 2012.

\bibitem{Hild14a}
Roland Hildebrand.
\newblock Minimal zeros of copositive matrices.
\newblock Preprint, \url{http://arxiv.org/abs/1401.0134}, 2014.

\bibitem{Loew03}
Raphael Loewy and Bit-Shun Tam.
\newblock C{P} rank of completely positive matrices of order 5.
\newblock {\em Linear Algebra Appl.}, 363:161--176, 2003.

\bibitem{Shak13b}
Naomi Shaked-Monderer, Abraham Berman, Immanuel~M. Bomze, Florian Jarre, and
  Werner Schachinger.
\newblock New results on the cp rank and related properties of
  co(mpletely~)positive matrices.
\newblock {\em Linear Multilinear Algebra}, to appear. Also available
  at:~arxiv.org/abs/1305.0737, 2013.

\bibitem{Shak13a}
Naomi Shaked-Monderer, Immanuel~M. Bomze, Florian Jarre, and Werner
  Schachinger.
\newblock On the cp-rank and minimal cp factorizations of a completely positive
  matrix.
\newblock {\em SIAM J. Matrix Anal. Appl.}, 34(2):355--�368, 2013.

\end{thebibliography}

\end{document} 